\section*{Conclusions}

We provide a reproducible comparative protocol for layered sparse linear classifiers on a real microarray
benchmark and a synthetic benchmark with known signal structure. Under repeated cross-validation, the three
analyzed methods (\texttt{cpl}, \texttt{svm}, \texttt{logreg}) can now be compared not only in terms of
predictive behavior across neuron depths and training-quality weighted voting, but also in terms of how quickly
they recover truly informative features and how rapidly later neurons become noise-dominated. The presented
protocol supports transparent, repeatable comparisons in which neuron depth is treated as an object of
analysis, not merely as a route to a final voted classifier.

The implemented model type is deliberately sparse and linear. This choice is justified by the small-sample,
high-dimensional setting: sparse formal neurons keep the selected feature sets inspectable, while the
complex-layer construction makes it possible to study whether successive sparse neurons expose complementary
signal after earlier selected features are removed. The comparison with sparse \texttt{svm} and
\texttt{logreg} baselines shows that this behavior depends on the loss function and the induced feature-count
regime, not only on the voting mechanism.

The main limitations are the small size of the real microarray benchmark, the controlled but necessarily
simplified synthetic data-generation process, and the use of calibrated analysis regimes rather than nested
hyperparameter selection. Future work should repeat the protocol on additional high-dimensional biological
datasets, study the stability of neuron-wise selected features, and develop stopping rules for identifying
when later neurons have become primarily noise-driven.
